\section{The state-dependent shift is onto}\label{sec:shift}

The weights of Lemma~\ref{lem:coordinate-weights} now provide the
map required by the absorption criterion. We first invert arbitrary
frozen weights, then bound their variation, and finally solve every
target by a finite-dimensional fixed-point argument.

\subsection{Frozen and state-dependent shifts}

First ignore the dependence on \(s\). The identity \(wv=1\) supplies
an inverse for any independently frozen choice of the pairs.

\begin{lemma}[Frozen shift]\label{lem:frozen}
Choose for every \((i,n)\) a pair \(v^{i,n},w^{i,n}\) satisfying
(W1), independently between blocks. Define \(S:D\to U\) by
\begin{equation}
 a_{i,0}=t_i,\qquad a_{i,n}=w^{i,n}x^{i,n},\qquad
 (Sx)^{i,n}=x^{i,n}+v^{i,n}(a_{i,n-1}-a_{i,n}).               \label{eq:frozen-shift}
\end{equation}
Then \(S\) is a bounded linear bijection with
\(\|S\|\leq5\) and \(\|S^{-1}\|\leq5\).
\end{lemma}

\begin{proof}
The row bounds give
\(\lvert a_{i,n-1}-a_{i,n}\rvert\leq4r_{i,n}(x)\);
for \(n=1\), the predecessor is just \(t_i\), and the same bound
holds. Thus
\[
                    \|(Sx)^{i,n}\|_\infty\leq5r_{i,n}(x).
\]
By \eqref{eq:radii-tail}, the block norms tend to zero, so \(Sx\in U\) and
\(\|S\|\leq5\).

For \(y\in U\), put \(b_{i,n}=w^{i,n}y^{i,n}\) and define
\begin{equation}
 t_i=b_{i,1},\qquad
 x^{i,n}=y^{i,n}-v^{i,n}(b_{i,n}-b_{i,n+1}).                  \label{eq:frozen-inverse}
\end{equation}
The block norms of \(y\) tend to zero off finite sets: otherwise
infinitely many coordinates of \(y\) would stay away from zero.
Since \(|b_{i,n}|\leq2\|y^{i,n}\|_\infty\), the scalar family
\((b_{i,n})\) belongs to \(c_0(\N^2)\). Formula \eqref{eq:frozen-inverse} therefore gives
\(t\in E\) and block norms of \(x\) tending to zero. Because the
blocks are finite, this is \(x\in D\). It also gives
\(\|x\|_D\leq5\|y\|_U\).

Applying \(w^{i,n}\) to \eqref{eq:frozen-inverse} yields
\(w^{i,n}x^{i,n}=b_{i,n+1}\), so \eqref{eq:frozen-shift} gives \(Sx=y\).
Conversely, when \(y=Sx\), applying \(w^{i,n}\) to \eqref{eq:frozen-shift} gives
\(b_{i,n}=a_{i,n-1}\); \eqref{eq:frozen-inverse} then recovers \(t_i\) and \(x^{i,n}\).
The formulas are mutual inverses and prove the inverse bound.
\end{proof}

For \(s\in B\), let \(S_s\) be the frozen shift using
\(v^\lambda(s),w^\lambda(s)\) in every block. Define
\begin{equation}
 g(0)=0,\qquad
 g(x)=S_{s_x}x,\quad s_x=\frac{x}{\|x\|_D}\quad(x\ne0).        \label{eq:g-definition}
\end{equation}
This map is positively homogeneous. Its weights are recomputed from
the input, so Lemma~\ref{lem:frozen} alone does not invert \(g\).

\begin{proposition}[Global distance bounds]\label{prop:global}
For every \(x,x'\in D\),
\begin{equation}
 \left(\frac15-24\eta\right)\|x-x'\|_D
 \leq\|g(x)-g(x')\|_U
 \leq(5+24\eta)\|x-x'\|_D.                                  \label{eq:global-bounds}
\end{equation}
In particular, \(g\) is injective and continuous.
\end{proposition}

\begin{proof}
Suppose first that \(x,x'\ne0\). Put
\(d=\|x-x'\|_D\) and \(M=\max\{\|x\|_D,\|x'\|_D\}\).
If \(\|x\|_D=M\), splitting
\(x/M-x'/\|x'\|\) at \(x'/M\), and using
\(|\|x\|-\|x'\||\leq d\), gives
\begin{equation}
                       \|s_x-s_{x'}\|_D\leq\frac{2d}{M}.     \label{eq:normalization}
\end{equation}
The other order is the same.

At each level \(\lambda\), freeze the entire pair
\((\bar v^\lambda,\bar w^\lambda)\) at whichever of \(s_x,s_{x'}\)
has larger \(r_\lambda\), breaking a tie arbitrarily. These choices
may vary with \(\lambda\), but Lemma~\ref{lem:frozen} still applies
to the resulting \(\bar S\). For \(e=x\) or \(x'\), the chosen pair
either agrees with the weights at \(s_e\) or the radius at \(s_e\)
is the smaller one. Thus (W2) and \eqref{eq:normalization} give
\begin{equation}
 r_\lambda(e)
 \max\{\|v^\lambda(s_e)-\bar v^\lambda\|_\infty,
       \|w^\lambda(s_e)-\bar w^\lambda\|_1\}
 \leq2\eta d.                                                \label{eq:frozen-weight-error}
\end{equation}
This remains valid when \(r_\lambda(e)=0\).

In block \(\lambda=(i,n)\), first change its vector \(v^\lambda\)
while keeping the two scalars computed with the original rows.
Their difference has modulus at most \(4r_\lambda(e)\), so \eqref{eq:frozen-weight-error}
bounds this change by \(8\eta d\). Changing the current row changes
its scalar by at most \(2\eta d\), since
\(\|e^\lambda\|_\infty\leq r_\lambda(e)\). If \(n>1\), changing the
predecessor row costs another \(2\eta d\), using \eqref{eq:frozen-weight-error} at \((i,n-1)\)
and \(\|e^{i,n-1}\|_\infty\leq r_{i,n-1}(e)\). For \(n=1\), the
predecessor \(t_i\) does not change. Since
\(\|\bar v^\lambda\|_\infty\leq1\), we have
\begin{equation}
              \|(S_{s_e}e-\bar S e)^\lambda\|_\infty
                    \leq12\eta d.                            \label{eq:frozen-block-error}
\end{equation}
Taking the supremum over blocks and adding the two endpoint errors,
\[
       \|g(x)-g(x')-\bar S(x-x')\|_U\leq24\eta d.
\]
The two frozen bounds give
\(\frac15d\leq\|\bar S(x-x')\|_U\leq5d\), proving \eqref{eq:global-bounds}.
If one endpoint is zero, Lemma~\ref{lem:frozen} applied to
\(S_{s_x}x\) gives the stronger constants \(1/5\) and \(5\).
Finally, \(1/5-24\eta>0\) for \(\eta=10^{-3}\).
\end{proof}

\subsection{Terminal scalars and Brouwer surjectivity}

The lower bound in \eqref{eq:global-bounds} proves injectivity and will prove that the
image is closed. It does not prove that every target has a preimage.
For that, we solve finite-coordinate targets with an extra terminal
scalar in each row.

\begin{proposition}[Onto-ness of the nonlinear shift]\label{prop:onto}
The map \(g:D\to U\) in \eqref{eq:g-definition} is onto. Together with
Proposition~\ref{prop:global}, it is a bi-Lipschitz bijection.
\end{proposition}

\begin{proof}
For \(N\geq1\), let \(U_N\) consist of the blocks with
\(1\leq i,n\leq N\), all other coordinates zero. Let \(D_N\) consist
of these blocks and \(t_1,\ldots,t_N\), all other input coordinates
zero. They are finite-dimensional, and
\begin{equation}
 \dim D_N=N+\sum_{i,n\leq N}|\Gamma_{i,n}|
          =\dim(U_N\oplus_\infty\R^N).                       \label{eq:finite-dimensions}
\end{equation}
For \(s\in B\), set
\begin{equation}
 S_s^{(N)}:D_N\to U_N\oplus_\infty\R^N,\qquad
 S_s^{(N)}x=
 \left(\bigl((S_sx)^{i,n}\bigr)_{i,n\leq N},
       \bigl(w^{i,N}(s)x^{i,N}\bigr)_{i\leq N}\right).        \label{eq:finite-shift}
\end{equation}
The last \(N\) coordinates retain the terminal scalars \(a_{i,N}\).

For clarity, this finite map has an explicit inverse. Given
\((y,(c_i)_{i\leq N})\), define
\[
 b_{i,n}=w^{i,n}(s)y^{i,n},\qquad
 a_{i,0}=b_{i,1},\qquad
 a_{i,n}=b_{i,n+1}\ (n<N),\qquad a_{i,N}=c_i,
\]
then put
\begin{equation}
 t_i=a_{i,0},\qquad
 x^{i,n}=y^{i,n}
          -v^{i,n}(s)(a_{i,n-1}-a_{i,n})\quad(n\leq N).      \label{eq:finite-inverse}
\end{equation}
Because \(a_{i,n-1}=b_{i,n}\), applying \(w^{i,n}(s)\) to \eqref{eq:finite-inverse}
gives \(w^{i,n}(s)x^{i,n}=a_{i,n}\). Substitution in \eqref{eq:finite-shift} returns
\((y,c)\). Conversely, applying the rows to \eqref{eq:finite-shift} recovers the
\(b_{i,n}\), and \eqref{eq:finite-inverse} recovers every input coordinate. Thus \eqref{eq:finite-inverse}
is the inverse.

The same estimates as for the frozen shift give
\begin{equation}
                  \|S_s^{(N)}\|\leq5,\qquad
                  \|(S_s^{(N)})^{-1}\|\leq5.                \label{eq:finite-bounds}
\end{equation}
For the inverse bound, if \(h_0=\|(y,c)\|\), then
\(|b_{i,n}|\leq2h_0\), all \(a_{i,n}\) have modulus at most \(2h_0\),
and \eqref{eq:finite-inverse} has block norm at most \(h_0+4h_0\). For the forward bound,
the block estimate is \(5\|x\|_D\) and the terminal coordinates
are at most \(2\|x\|_D\). The finitely many weights in \eqref{eq:finite-inverse} are
continuous in \(s\). Hence \(s\mapsto(S_s^{(N)})^{-1}\) is continuous
as a map into the linear operators between these finite-dimensional
spaces.

Figure~\ref{fig:terminal} displays the terminal coordinate and the
single possible outgoing block of a finite row.
\begin{figure}[htbp]
\centering
\begin{tikzpicture}[
  >=Latex,
  every node/.style={font=\small},
  block/.style={draw,minimum width=29mm,minimum height=10mm,align=center},
  scalar/.style={draw,minimum width=21mm,minimum height=8mm,align=center}
]
\node[scalar] (a0) at (0,0) {\(a_{i,0}=t_i\)};
\node[block] (y1) at (3.0,0) {\(y^{i,1}\)\\\(w^{i,1}y^{i,1}=a_{i,0}\)};
\node (dots) at (5.2,0) {\(\cdots\)};
\node[block] (yn) at (7.4,0) {\(y^{i,N}\)\\\(w^{i,N}y^{i,N}=a_{i,N-1}\)};
\node[scalar] (terminal) at (10.7,0) {\(c_i=a_{i,N}\)\\extra target};
\draw[->] (a0) -- (y1);
\draw[->] (y1) -- (dots);
\draw[->] (dots) -- (yn);
\node[block] (spill) at (10.7,-1.8)
 {\((i,N+1)\) output\\\(v^{i,N+1}c_i\)};
\draw[->] (terminal) -- node[right,align=left]
 {\(x^{i,N+1}=0\)} (spill);
\node[align=left,anchor=west] at (0,-1.8)
 {Finite target: \((y,c)\).\\
  Brouwer is applied with \(c=0\).};
\end{tikzpicture}
\caption{A row of length \(N\) with frozen weights and independent
extra target scalar \(c_i=a_{i,N}\). The block readouts recover the
predecessor scalars, while \(c_i\) records the scalar not read by any
block in the finite row. At the nonlinear
fixed point it is set to zero, which kills the only possible output
spillover into the next block.}
\label{fig:terminal}
\end{figure}


Let \(0\ne y\in U_N\), and set \(h=\|y\|_U>0\). On the closed ball
\(\{x\in D_N:\|x\|_D\leq5h\}\), define
\begin{equation}
 s(x)=\frac{x}{\max\{\|x\|_D,h/6\}},\qquad
 T(x)=(S_{s(x)}^{(N)})^{-1}(y,0).                            \label{eq:brouwer-map}
\end{equation}
The denominator is positive even at \(x=0\), so \(s(x)\in B\)
depends continuously on \(x\). Equations \eqref{eq:finite-inverse}--\eqref{eq:finite-bounds} show that \(T\)
is a continuous self-map of this finite-dimensional closed ball.
Brouwer's fixed-point theorem gives \(x=T(x)\). At that fixed point,
\[
       h=\|(y,0)\|=\|S_{s(x)}^{(N)}x\|\leq5\|x\|_D,
\]
and therefore \(\|x\|_D\geq h/5>h/6\). The state in \eqref{eq:brouwer-map} is now
the normalized state \(s_x\), not the regularized state used near zero.
The blocks of \(g(x)\) in the \(N\)-square are \(y\), and \eqref{eq:finite-shift} gives
\(a_{i,N}=0\) for every \(i\leq N\).

It remains to check the blocks omitted from the finite equation.
The input is zero outside \(D_N\). The only outside block with a
possibly nonzero predecessor scalar is \((i,N+1)\) for \(i\leq N\).
There
\[
 (g(x))^{i,N+1}
       =v^{i,N+1}(s_x)a_{i,N}=0.
\]
All other outside blocks have zero input and zero predecessor scalar.
Thus \(g(x)=y\) in \(U\). The zero target is \(g(0)\), so
\(U_N\subset g(D)\) for every \(N\).

Every finitely supported element of \(U\) belongs to some \(U_N\),
and such elements are dense. The image \(g(D)\) is closed: if
\(g(x_k)\) converges in \(U\), the positive lower bound \eqref{eq:global-bounds} makes
\((x_k)\) Cauchy in the Banach space \(D\); its limit \(x\) satisfies
\(g(x_k)\to g(x)\) by continuity. Hence \(g(D)\) is both dense and
closed in \(U\), and is all of \(U\).
\end{proof}

Let \(P:D\to U\) be the projection and, for \(s\in B\), put
\begin{equation}
                         h(s)=(S_s-P)s.                      \label{eq:h-definition}
\end{equation}
The frozen lemma ensures \(h(s)\in U\), and \eqref{eq:correction-coordinate} is exactly its
\(\gamma\)-coordinate. Notice that \(S_s s\) uses \(s\) itself
as state even when \(\|s\|<1\), whereas \(g(s)\) uses \(s/\|s\|\).
The lifting argument will use \eqref{eq:h-definition} on \(B\) and \eqref{eq:g-definition} only on
normalized states. By \eqref{eq:correction-bounds}, the correction coordinates have common
bounds
\begin{equation}
 A:=\sup_\gamma\|f_\gamma\|_\infty\leq4,\qquad
 L_0:=\sup_\gamma\Lip(f_\gamma)\leq4+12\eta.                  \label{eq:base-bounds}
\end{equation}
